% GENERATED FILE. Edit canonical lessons, then run npm run generate:books.
% canonical-source-hash: fnv1a64:368581c013556efc
% canonical-lessons: MR-C36-jevha, MR-C36-lyaavar, MR-C36-sudhaa, MR-R36-joined

\chapter{When, After, and Too}
\label{ch:mr-time-joins}
\hlchaptermodality{36}

\begin{chapteropening}
\textbf{By the end of this chapter:} \emph{I can place two events in time with jevhaa ... tevhaa or the shorter -lyaavar, add something with sudhaa, and run every joining word in the book once.}

Complete MR-R36-joined and say five sentences with a clause underneath -- a reason, a content, an aim, a same-time and an after -- naming the job of each.
\end{chapteropening}

\section[jevhā … tevhā …]{\mr{जेव्हा} … \mr{तेव्हा} … (jevhā … tevhā …) — when, then}

\label{lesson:MR-C36-jevha}



\begin{quote}
Retrieve the two items before this one by name: \textbf{\mr{की}}, which hangs a whole sentence on a thinking verb, and \textbf{-\mr{ण्यासाठी}}, which turns a verb into an aim. Both put a clause underneath. This one does it in time.
\end{quote}

\begin{grammarlens}[title={Grammar Lens: the two-word frame}]
\begin{quote}
\textbf{\mr{जेव्हा} \mr{मी} \mr{वाचतो} \mr{तेव्हा} \mr{मला} \mr{समजतं}.} — \emph{jevhā mī vācto tevhā malā samajtaṁ.} — \textbf{When I read, then I understand.}
\end{quote}

Marathi does not leave \emph{when} on its own. It opens with \textbf{\mr{जेव्हा}} and answers with \textbf{\mr{तेव्हा}} in the second half, and the second word is not optional padding — it is what tells you the first clause has finished.

You have met this machinery before without being told its name. \textbf{\mr{जेव्हा}} is built on the relative stem \textbf{\mr{ज}-}, \textbf{\mr{तेव्हा}} on the demonstrative \textbf{\mr{त}-}, and the same two stems sit under \textbf{\mr{जो}/\mr{तो}}. English keeps the pair in \emph{when … then} and drops the \emph{then}; Latin kept both in \emph{cum … tum}. Marathi keeps both, always.

Because the frame is spoken in the order the events happen, it is also the easiest long sentence in the book to say: first clause, \textbf{\mr{तेव्हा}}, second clause.
\end{grammarlens}

\subsection*{Your turn}
Say these aloud:

\begin{itemize}
\raggedright
  \item \emph{jevhā mī vācto tevhā malā samajtaṁ}
  \item the pair on its own — \emph{jevhā … tevhā}
  \item which of the two stems is the relative one — \textbf{\mr{ज}-}
\end{itemize}

\begin{tcolorbox}[breakable,colback=teal!4,colframe=teal!35!black,title={Before you move on}]
What must follow a \textbf{\mr{जेव्हा}} clause? (\textbf{\mr{तेव्हा}}.) What does that second word do for the listener? (\textbf{Tells them the first clause has ended}.)

Sources: \href{https://en.wiktionary.org/wiki/\%E0\%A4\%9C\%E0\%A5\%87\%E0\%A4\%B5\%E0\%A5\%8D\%E0\%A4\%B9\%E0\%A4\%BE}{Wiktionary: \mr{जेव्हा}}.
\end{tcolorbox}
\section[-lyāvar]{-\mr{ल्यावर} (-lyāvar) — after doing}

\label{lesson:MR-C36-lyaavar}



\begin{quote}
Say the two-clause frame you got last lesson: \textbf{\mr{जेव्हा} \mr{मी} \mr{वाचतो} \mr{तेव्हा} \mr{मला} \mr{समजतं}}. Nine syllables of scaffolding for two events. Marathi has a shorter way.
\end{quote}

\begin{grammarlens}[title={Grammar Lens: one ending instead of two clauses}]
\begin{quote}
\textbf{\mr{वाचल्यावर} \mr{मला} \mr{समजतं}.} — \emph{vācalyāvar malā samajtaṁ.} — \textbf{After reading, I understand.}
\end{quote}

The swap is the same shape as \textbf{-\mr{ण्यासाठी}}: take \textbf{-\mr{णे}} off the verb and put \textbf{-\mr{ल्यावर}} on. The result is one word that means \emph{after doing it}, and it stands at the front of the sentence with no second half required.

\noindent
\begin{tabularx}{\linewidth}{@{}>{\raggedright\arraybackslash}X>{\raggedright\arraybackslash}X@{}}
\toprule
\textbf{dictionary form} & \textbf{after-doing form} \\
\midrule
\mr{वाचणे} & \mr{वाचल्यावर} \\
\mr{शिकणे} & \mr{शिकल्यावर} \\
\bottomrule
\end{tabularx}

Taken apart, \textbf{-\mr{ल्यावर}} is a past participle in \textbf{-\mr{ल}-} with \textbf{\mr{वर}} (\textquotedblleft{}on, upon\textquotedblright{}) behind it. \emph{After reading} is literally \emph{upon having-read}. English reaches for the same preposition in \emph{upon reading}, which makes this one of the few Marathi endings you can feel your way into rather than memorise.

\textbf{\mr{जेव्हा}} says the two things happen at once; \textbf{-\mr{ल्यावर}} says one finished before the other began. Choose by which you mean, not by which is shorter.
\end{grammarlens}

\subsection*{Your turn}
Say these aloud:

\begin{itemize}
\raggedright
  \item \emph{vācṇe} $\to$ \emph{vācalyāvar}, then \emph{shikṇe} $\to$ \emph{shikalyāvar}
  \item \emph{vācalyāvar malā samajtaṁ}
  \item which of the two says \emph{at the same time} — \emph{jevhā}
\end{itemize}

\begin{tcolorbox}[breakable,colback=teal!4,colframe=teal!35!black,title={Before you move on}]
What does \textbf{-\mr{ल्यावर}} replace? (\textbf{A whole \mr{जेव्हा} … \mr{तेव्हा} frame}.) What is the \textbf{\mr{वर}} in it doing? (\textbf{It means \emph{upon} — upon having done it}.)
\end{tcolorbox}
\section[sudhā]{\mr{सुद्धा} (sudhā) — too, and \textquotedblleft{}not even\textquotedblright{}}

\label{lesson:MR-C36-sudhaa}



\begin{quote}
\textbf{\mr{आणि}} joins two things in one sentence. But often the first thing was said a minute ago, by somebody else. There is no gap to put \textbf{\mr{आणि}} in.
\end{quote}

\subsection*{You'll want to know: \mr{सुद्धा}}
\begin{quote}
\textbf{\mr{मला} \mr{दूध} \mr{सुद्धा} \mr{आवडतं}.} — \emph{malā dūdh sudhā āvaḍtaṁ.} — \textbf{I like milk too.}
\end{quote}

\textbf{\mr{सुद्धा}} goes straight after the word it adds — not between two things, but behind one. That is what lets it reach back to something already said. It is the joining word for a CONVERSATION rather than for a sentence.

\begin{grammarlens}[title={Grammar Lens: the same word in front of a denial}]
\begin{quote}
\textbf{\mr{मला} \mr{दूध} \mr{सुद्धा} \mr{आवडत} \mr{नाही}.} — \emph{malā dūdh sudhā āvaḍat nāhī.} — \textbf{I don't even like milk.}
\end{quote}

Nothing about \textbf{\mr{सुद्धा}} changed. The sentence around it turned negative, and \emph{too} became \emph{not even} on its own. English does the identical thing: \emph{I like it too} against \emph{I don't even like it}. One word, two readings, and the negation picks which.

\textbf{\mr{सुद्धा}} is Sanskrit \textbf{\mr{शुद्ध}} (\emph{śuddha}), \textquotedblleft{}pure.\textquotedblright{} \emph{Purely that} hardened into \emph{even that}, which is the ordinary road a purity word travels.
\end{grammarlens}

\subsection*{Your turn}
Say these aloud:

\begin{itemize}
\raggedright
  \item \emph{malā chahā āvaḍto} — then \emph{malā dūdh sudhā āvaḍtaṁ}
  \item \emph{malā dūdh sudhā āvaḍat nāhī}
  \item where \textbf{\mr{सुद्धा}} stands — \textbf{right after the word it adds}
\end{itemize}

\begin{tcolorbox}[breakable,colback=teal!4,colframe=teal!35!black,title={Before you move on}]
Where does \textbf{\mr{सुद्धा}} stand? (\textbf{Directly after the word it adds}.) What turns \emph{too} into \emph{not even}? (\textbf{Negating the verb; \mr{सुद्धा} itself does not change}.)

Sources: \href{https://en.wiktionary.org/wiki/\%E0\%A4\%B8\%E0\%A5\%81\%E0\%A4\%A6\%E0\%A5\%8D\%E0\%A4\%A7\%E0\%A4\%BE}{Wiktionary: \mr{सुद्धा}}.
\end{tcolorbox}
\section[āṇi · kiṁvā · paṇ · kāraṇ · kī · jevhā]{Everything that joins}

\label{lesson:MR-R36-joined}



\begin{quote}
Seven chapters ago you could name a hundred Marathi things and say nothing about two of them at once. Name the six joining words in this lesson's title without looking.
\end{quote}

\subsection*{Your turn — flat joins}
Say these aloud:

\begin{itemize}
\raggedright
  \item \emph{chahā āṇi dūdh} · \emph{chahā kiṁvā dūdh} · \emph{chahā āhe paṇ dūdh nāhī}
  \item \emph{nā chahā nā dūdh} · \emph{ek chahā āṇi dusrā pāṇī}
  \item \emph{malā dūdh sudhā āvaḍtaṁ} · \emph{malā dūdh sudhā āvaḍat nāhī}
  \item \emph{chahā āhe nā?} · \emph{chahā? — nako, dhanyavād}
\end{itemize}

\subsection*{Your turn — clauses underneath}
Say each of the five, and after each one name what the underneath clause is doing: a reason, a content, an aim, a same-time, an after.

\begin{tcolorbox}[breakable,colback=teal!4,colframe=teal!35!black,title={Before you move on}]
Which joining words treat the two halves as EQUALS? (\textbf{\mr{आणि}, \mr{किंवा}, \mr{पण}, \mr{ना} … \mr{ना}, \mr{एक} … \mr{दुसरा}}.) Which ones put one half UNDERNEATH the other? (\textbf{\mr{कारण}, \mr{की}, -\mr{ण्यासाठी}, \mr{जेव्हा} … \mr{तेव्हा}, -\mr{ल्यावर}}.) And which single syllable both asks a yes-or-no question and asks \emph{why}, depending only on where it stands? (\textbf{\mr{का}}.) When it is dropped altogether, what carries the question instead? (\textbf{A rising voice} — and on paper, the single \textbf{?}, with \textbf{!} its partner.) Then stop.
\end{tcolorbox}
